% page 75 du fac-simile (image p-074 du scan)
% bloc 162.6 x 229.0 mm ; marge gauche 39.4 mm
\pgc{17.780mm}{0.000mm}{
\l{\cel{55}67}
\l{}
\l{}
\l{\sou{bitumo}. Substanco minerala, liquida -nafto) o solida (asfalto),}
\l{\cel{3}bruna. - DEFIRS.}
\l{}
\l{\sou{bivakar}. (\sou{netrans.}) Kampar sen tendi, havante cielo kom plafono,}
\l{\cel{3}generale cirkum fairi), nur dum un nokto. - DEFIRS.}
\l{}
\l{\sou{bizantina}. Qua prizas disipar sua tempo en diskutar subtilaji}
\l{\cel{3}ne-utila. - DEFIRS.}
\l{}
\l{\sou{bizelo.}\cel{2}Bordo (di planko, cizelilo, e c.) tranchita oblique.}
\l{\cel{3}- FS.}
\l{}
\l{\sou{bizono}. (\sou{zool}.)Bovo sovaja di Usa, diqua la dorso esas giboza,}
\l{\cel{3}la korni kurta e rondatra, kun longa barbo. - L.\cel{1}\sou{bos}}
\l{\cel{3}\sou{americanus}. - DEFIRS.}
\l{}
\l{\sou{blamar}.\cel{2}(\sou{trans., pro, pri}). Desaprobar ulu quan onu judikas}
\l{\cel{3}kom aginta male. - DEFIS.}
\l{}
\l{\sou{blanka}. Di qua la koloro esas pasable simila a lakto, a}
\l{\cel{3}nivo-floki, e c. - DEFIRS.}
\l{}
\l{\sou{blasfemar}. (netrans.) Pronuncar vorti qui ofensas la deajo.-}
\l{\cel{3}DEFIS.}
\l{}
\l{\sou{blato}. (\sou{zool.}) Insekto ortoptera bruna, qua odoras fetide,}
\l{\cel{3}ajila, devorema, qua devoras la nutrivo-provizuri, rodas}
\l{\cel{3}la stofi, e c. - L.\cel{1}\sou{blatta}. - DFI.}
\l{}
\l{\sou{blazar}. (\sou{trans}.) Obtuzigar per impresi granda ed ofte iterata.}
\l{\cel{3}- DEF.}
\l{}
\l{\sou{blazono}. Ensemblo de emblemi konsiderata kom signo distingiva}
\l{\cel{3}di familio nobela, di urbo, di korporaciono, e c., quan}
\l{\cel{3}onu reprezentis sur la skudi, la baneri, la siglili, e c.}
\l{\cel{3}- DEFIS.}
\l{}
\l{\sou{blendo}. (\sou{kemio}).Erco de zinko-sulfuro, ordinare unionita}
\l{\cel{3}a la mineralo-strati de plombo, de arjento. - DEFIRS.}
\l{}
\l{\sou{blezar}. (\sou{netrans}.) Pronuncar \textquotedbl{}s\textquotedbl{} quale \textquotedbl{}th\textquotedbl{} Angla o \textquotedbl{}c\textquotedbl{}}
\l{\cel{3}(avan \textquotedbl{}a\textquotedbl{} e \textquotedbl{}i\textquotedbl{}) Hispana (t.e. kun la lango-pinto inter la}
\l{\cel{3}denti.) - F.}
\l{}
\l{\sou{blinda}. Sen vidiveso. - DE.}
\l{}
\l{\sou{bloko}. La toto di korpo ; amaso granda e pezoza. - DEFIRS.}
\l{}
\l{\sou{blokhauso}. (\sou{milit.}) Fuorteto regiono-sola, ek ligno, ordinare}
\l{\cel{3}juntita a fortifikajo precipua per koridoro subsula. -DEF.}
}
